\documentclass[11pt]{article}
\usepackage[margin=0.9in]{geometry}
\usepackage{amsmath,amssymb,amsthm,booktabs,mathrsfs}
\usepackage[T1]{fontenc}
\usepackage{lmodern,microtype}
\usepackage[colorlinks=true,linkcolor=blue,urlcolor=blue]{hyperref}
\newcommand{\norm}[1]{\left\lVert#1\right\rVert}
\newtheorem{theorem}{Theorem}[section]
\newtheorem{lemma}[theorem]{Lemma}
\newtheorem{proposition}[theorem]{Proposition}
\newtheorem{corollary}[theorem]{Corollary}
\theoremstyle{remark}
\newtheorem{remark}[theorem]{Remark}
\title{Toward a four-dimensional Yang--Mills construction\\
\large Notebook 38: positive full-measure embeddings and local defect bounds}
\author{Research proof draft}
\date{17 October 2026}
\setlength{\emergencystretch}{5em}
\begin{document}
\maketitle
\begin{abstract}
We move beyond the homogeneous local integral to exact statements
about finite positive Wilson measures. A specified measurable
completion of the actual local coarse map admits positive linear
good and remainder operators whose sum is the exact pushforward.
Their iterates retain every history and the full normalization.
Under literal input-domain containment and the established local
Ward identities, the good measure has the actual local analytic
germ on each fixed finite regulator.

Independently, a site-plane slab argument bounds prescribed
large-curvature cells under the normalized Wilson measure, uniformly
in the total volume. A compact-group conditional Laplace theorem
controls fluctuations in a fixed box with small boundary data.
A boundary-only tree gauge joins these results: prescribed
separated boxes have a joint exponentially small probability of
failure of explicit local gauge/minimizer windows. These are
estimates on the entire bare Wilson measure, with the denominator
included.

The local box windows have not been identified with the averaged
constraint windows of the iterated renormalization map. The
all-history decomposition has not been bounded in a multiscale
complete-density norm. We record this remaining interface, a
precise match to the published large-field induction, and a
regularity counterexample ruling out one insufficient shortcut.
The continuum existence theorem is not established.
\end{abstract}

\section*{Relation to the construction and independence of inputs}

Notebook 37 establishes a positive logarithmic drift and arbitrarily
long finite weak homogeneous local histories within the working
proof chain and its explicitly stated upstream hypotheses. Those
histories alone do not represent the complete Wilson measure.
Here we address that next interface.

There are two distinct kinds of result below. The positive
embedding uses the previously defined ordinary composite coarse
map, stationary chart, pivot measure and local gauge Ward identity.
Its local-germ assertion is conditional on literal containment of
the real integration window in the represented input density.
By contrast, the bare Wilson probability results are proved
directly from compact-group integration, finite lattice geometry,
and positivity. They do not require the earlier multiscale
determinant calculation.

All field and coupling arguments in the new probability bounds
are real. The compact gauge groups and link Haar measures are
normalized to mass one. Unless explicitly stated otherwise, a
fixed-box constant may depend on its side length and the group,
but not on the surrounding torus volume. Such a constant is not
claimed uniform in a growing physical box at vanishing lattice
spacing. The general \(SU(N)\) bare-measure lemmas do not extend
all earlier \(SU(2)\) construction inputs to arbitrary groups.

The source repository remains pinned at
\texttt{e392ef4501\allowbreak 3a581664ba\allowbreak b76b908671\allowbreak 70b30fe10e}. This document is
a mathematical research draft with separate internal audits, not
an assertion of external verification of the preceding notebooks.
The exact status of each result and each missing implication is
part of its statement.



% No preamble. Prefix fps-. 17 October 2026.
\section{An exact positive finite-measure embedding of the retained sector}
\label{fps-section}

Fix a finite compatible periodic lattice and a compact connected matrix
gauge group, with normalized link Haar measure. This section constructs
an exact positive decomposition of the original finite Wilson measure
under a specified real measurable coarse map. It does not estimate the
complementary histories or identify the full pushforward with the
homogeneous analytic density.

The primitive ordinary mean, composite ratio \(R=\ell^s\), direct-\(R\)
internal contour gauge, physical centered branch \(V_*(W)\), and retained
pivot chart are the ones already established for the local macrostep.
A global completion is fixed explicitly below. It agrees with that
analytic map near the identity, but is not identified with an unspecified
different large-field completion in the deposited account.

\subsection{The global real map and its exact compact-gauge normalization}
\label{fps-completion-section}

For a contour tuple \((U_1,\ldots,U_n)\) from \(y\) to \(x\), set
\(T=U_1\). On a sufficiently small conjugation-invariant identity
neighborhood, let \(\Phi\) be the actual selected local mean on the
reference-normalized tuple, with its original ordering, multiplicities
and reference. On the Borel good set where all \(U_jT^{-1}\) lie in that
neighborhood, define
\[
 P(U_1,\ldots,U_n)=\Phi(U_1T^{-1},\ldots,U_nT^{-1})T.
\]
On its complement put \(P=T\). The good set is endpoint-invariant, and
local conjugation covariance gives, everywhere on the real group,
\[
 P(gU_1h^{-1},\ldots,gU_nh^{-1})=gP(U_1,\ldots,U_n)h^{-1}.
\]
Right covariance proves agreement with the original small-field mean.
No alternative group mean is substituted.

For each primitive ordinary coarse bond, compute these completed
internal contours and its original center string \(T_c\). If every
transported relative string \(U_xT_c^{-1}\) lies in the specified small
logarithm domain, use the original formula
\[
 \exp\left[\ell^{-d}\sum_x\log(U_xT_c^{-1})\right]T_c;
\]
otherwise use \(T_c\). This defines a Borel gauge-covariant primitive
map. Compose these primitive completions to obtain \(F\). Complete the
independent direct-\(R\) internal contour \(P_R\) in the same way.
Thus, with \(r\) restriction to final block centers,
\begin{equation}
 F(U^v)=F(U)^{r(v)},\qquad
 P_R(U^v;y,x)=v(y)P_R(U;y,x)v(x)^{-1}.
 \label{fps-global-covariance}
\end{equation}
The completion need not be continuous at its selected boundaries and
is not asserted globally holomorphic. The original fine Wilson measure
is unchanged: \(F\) is now a fully specified coarse observable of it.
No equality with a direct-\(R\) ordinary average is used.

Let \(H_R\) be the site-gauge group equal to the identity at every final
block center, with product normalized Haar, and let \(n_s\) count its
free site variables. For positive real \(g\), define
\begin{equation}
 q(h)=1-\operatorname{Re}\tau(h),\quad
 G_R(U)=\sum_{y,x\ne y}q(P_R(U;y,x)),\quad
 c(g)=\int_Ge^{-q(h)/g^2}\,dh>0.
 \label{fps-compact-normalizer}
\end{equation}
This is the exact compact integral, not a Gaussian approximation.
Equation~\eqref{fps-global-covariance} and independence of the site
Haar variables give, for every real configuration,
\begin{equation}
 \int_{H_R}e^{-G_R(U^u)/g^2}\,du=c(g)^{n_s}.
 \label{fps-orbit-normalization}
\end{equation}
Indeed its individual contour factors are
\(P_R(U;y,x)u(x)^{-1}\), each with Haar integral \(c(g)\).

\begin{lemma}[Pointwise orbit acceptance]
\label{fps-markov-lemma}
The kernel
\begin{equation}
 K_g(U,dV)=c(g)^{-n_s}\int_{H_R}
       e^{-G_R(U^u)/g^2}\,\delta_{U^u}(dV)\,du
 \label{fps-markov}
\end{equation}
is a probability kernel supported on \(F(V)=F(U)\).
For every Borel \(0\le\chi(V,W)\le1\), put
\(p_\chi(U)=\int K_g(U,dV)\chi(V,F(U))\). Then for every
finite positive measure \(\mu\),
\begin{equation}
 F_*\mu=F_*(p_\chi\mu)+F_*((1-p_\chi)\mu),\qquad
 0\le p_\chi\le1.
 \label{fps-markov-split}
\end{equation}
If \(\mu\) is \(H_R\)-invariant, then
\begin{equation}
 K_g^*\mu(dV)=c(g)^{-n_s}e^{-G_R(V)/g^2}\mu(dV).
 \label{fps-invariant-insertion}
\end{equation}
\end{lemma}
\begin{proof}
Normalization follows from \eqref{fps-orbit-normalization}; based gauge
invariance of \(F\) gives the support statement. The split is therefore
pointwise acceptance/rejection. For the last assertion, in the integral
against each test function change \(U\) to \(U^u\), using invariance of
\(\mu\), then integrate \(du\). For noninvariant input the corresponding
formula uses its \(H_R\)-average, not its original density.
\end{proof}

\subsection{The actual local density and its positive domination}
\label{fps-local-section}

Let \(\mu\) be invariant under the full current-site gauge group and
write its Haar Lebesgue decomposition
\(\mu=\rho\,m_E+\mu_{\rm sing}\). Both parts are invariant by uniqueness
of the decomposition and Haar invariance. Use only \(\rho m_E\) for the
local coarea operator; every remaining positive singular part will stay
in the complement.

For \(W\) in a small real domain \(D\), set \(V=V_*(W)\). The actual
chart \(\Phi_V(z)\) solves the relative constraint exactly and leaves
every retained coordinate equal to \(z\). Choose an injective local
sheet with a fixed real retained window \(\mathcal B_z\). An additional
factor \(0\le\chi_{\rm extra}\le1\) may impose any further original
pivot, contour or reconstructed-field restrictions; those factors are
not silently dropped.

Let \(n_e,n_c\) count positive current and coarse bonds, and set
\(n_z=(\dim\mathfrak g)(n_e-n_c)\). In real Lie coordinates orthonormal
for the chosen trace form, write
\(d(e^{iX})=h_0j(X)dX\), with \(j(0)=1\).
For the actual pivot blocks \(A_{c,p}\), define
\[
 D_R(V)=\prod_c\det(A_{c,p})^{-1}.
\]
All real small-chart orientations are fixed positive. Let
\(j_c^{\rm str}(V;K_Vz)\) be the logarithm of
\(\det A_{c,p}/\det D_pF_V(\Phi_V(z);c)\), zero at \(z=0\).
The exact local density relative to coarse product Haar is
\begin{align}
 \ell_\mu(W)
 &=c(g)^{-n_s}h_0^{n_e-n_c}D_R(V)
   \int_{\mathcal B_z}\rho(e^{i\Phi_V(z)}V)
       e^{-G_R(e^{i\Phi_V(z)}V)/g^2}\notag\\
 &\qquad\times\prod_b j(\Phi_V(z)(b))
       \exp\left[\sum_c j_c^{\rm str}(V;K_Vz)\right]
       \chi_{\rm extra}\,dz.
 \label{fps-local-density}
\end{align}
The target group delta contributes \(h_0^{-1}\) per coarse bond.
The square ambient straightening and subsequent pivot elimination give
this expression; a rectangular \(\det D\Phi\) is not substituted for
the constraint Jacobian. Equivalently its unrearranged pivot factor is
the product of inverse nonlinear pivot derivatives.
For measurable \(\rho\), the local joint change of variables and
Tonelli prove \eqref{fps-local-density} almost everywhere. It is a
nonnegative measurable function. Applying
\eqref{fps-invariant-insertion} to this actual local sheet gives
\begin{equation}
 \nu:=F_*\mu\ \ge\ 1_D\ell_\mu\,m_C.
 \label{fps-local-domination}
\end{equation}
No global density or regularity for \(\nu\) is required.

If the literal input chart has
\(\rho(U)=Z_{\rm in}\exp[-S_M(U)/g^2+E_{\rm raw}(U)]\),
put \(z=g\xi\). With the actual pulled-back action-plus-gauge Hessian
\(A_M(V)\), its exact cubic-and-higher remainder and all measure factors
in \(I\), the same formula is
\begin{align}
 \ell_\mu(W)
 &=Z_{\rm in}c(g)^{-n_s}h_0^{n_e-n_c}D_R(V)
       g^{n_z}(2\pi)^{n_z/2}\det A_M(V)^{-1/2}\notag\\
 &\quad\times e^{-S_M(V)/g^2+E_{\rm raw}(V)}
 \mathbb E_{A_M(V)^{-1}}
       [\chi_{\rm window}\chi_{\rm extra}e^{I(V,g,\xi)}].
 \label{fps-gaussian-density}
\end{align}
Here \(I\) contains \(-g^{-2}R_3\), all Haar and straightening
logarithms, and the exact inherited fluctuation difference. Its field
argument is \(\Phi_V(g\xi)\), not \(g\Phi_V(\xi)\).
This is an exact positive-real identity only when the input density
representation holds at every reconstructed field of the window.
Holomorphic continuation alone does not furnish that condition.

\subsection{An invariant positive good/remainder pair}
\label{fps-invariant-section}

Choose a fixed continuous \(\theta\ge0\) supported in \(D\), positive
near the identity. Let \(\mathcal G_C\) be the full coarse-site gauge
group with normalized Haar. Define
\[
 m_\theta(W)=\int_{\mathcal G_C}\theta(W^v)\,dv.
\]
It is invariant, and positive on the orbit saturation of
\(\{\theta>0\}\). With value zero when \(m_\theta=0\), set
\begin{equation}
 a_\mu(W)=\frac{\int_{\mathcal G_C}
                   \theta(W^v)\ell_\mu(W^v)\,dv}{m_\theta(W)},
 \qquad \mathcal A_g\mu=a_\mu m_C.
 \label{fps-orbit-extension}
\end{equation}
Only real values of the locally computed density are used, where
\(\theta\) is supported. This is not a globally analytic extension.

\begin{theorem}[Exact invariant positive split]
\label{fps-positive-split}
On the cone of full-gauge-invariant finite positive measures,
\begin{equation}
 0\le\mathcal A_g\mu\le F_*\mu,\qquad
 \mathcal R_g\mu:=F_*\mu-\mathcal A_g\mu\ge0.
 \label{fps-positive-operators}
\end{equation}
Both operators are linear, preserve gauge invariance, and satisfy
\(\mathcal A_g+\mathcal R_g=F_*\), including raw mass.
\end{theorem}
\begin{proof}
The full pushforward is invariant by
\eqref{fps-global-covariance}. Its Haar-absolutely-continuous part has
an invariant nonnegative density \(f\), chosen by compact gauge
averaging. Equation~\eqref{fps-local-domination} implies
\(f\ge\ell_\mu\) almost everywhere on \(D\). For almost every \(W,v\),
\[
 \theta(W^v)\ell_\mu(W^v)\le\theta(W^v)f(W^v)
                              =\theta(W^v)f(W).
\]
Integrating in \(v\) and dividing on \(m_\theta>0\) proves
\(a_\mu\le f\). The positive singular part remains in the remainder.
Tonelli, or truncation of the positive factor \(m_\theta^{-1}\),
justifies this near its zero set. Linearity follows from the linear
Haar-absolutely-continuous projection, local fiber integration and
fixed \(\theta\). The remainder is positive by the proved inequality.
\end{proof}

The proof cannot be replaced by averaging a measure inequality and
then dividing it by \(m_\theta\). It uses the invariant density bounding
each orbit value before normalization. No uniform lower bound on
\(m_\theta\) is needed. This remainder need not be the literal rejection
in \eqref{fps-markov-split}: the normalized output average redistributes
the local sector along coarse gauge orbits. Both splits are exact.

\subsection{The local germ is retained on a fixed finite regulator}
\label{fps-germ-section}

\begin{lemma}[Finite orbit-overlap control]
\label{fps-overlap}
Let \(\ell\) be an analytic local density on a product potential ball of
radius \(r\) on a connected finite coarse lattice. Assume its total
infinitesimal gauge Ward identity holds there and it is invariant under
constant color conjugations. There is \(\epsilon>0\), possibly depending
on the finite graph, such that
\begin{equation}
 W,W^v\in D_\epsilon\quad\Longrightarrow\quad
 \ell(W^v)=\ell(W).
 \label{fps-overlap-identity}
\end{equation}
\end{lemma}
\begin{proof}
Use a unitarily invariant matrix norm, paying the fixed color-coordinate
equivalence constants in \(r\). Choose a rooted spanning tree with
maximal root-path length \(D_{\rm tree}\). Write \(v_x=h_xc\), where
\(c=v_o\), \(h_o=I\). Constant conjugation does not increase the
invariant field norm. On each tree edge \(xy\),
\[
 h_y=(W^v_{xy})^{-1}h_xW^c_{xy}.
\]
Telescoping unitary products gives
\(\|h_x-I\|\le C D_{\rm tree}\epsilon\).
Choose \(\epsilon\) so every \(h_x=e^{i\lambda_x}\) uses one small
logarithm chart, with
\(\|\lambda_x\|\le C D_{\rm tree}\epsilon\).
Along \(h_x(t)=e^{it\lambda_x}\), the transformed angles are bounded by
\(C(1+D_{\rm tree})\epsilon\). Taking
\(\epsilon\le cr/(1+D_{\rm tree})\) keeps the full real path inside
the Ward domain. Integrate that identity, then apply constant
conjugation invariance. No global logarithm of the original \(v\)
was assumed.
\end{proof}

For the actual chosen sector, this application uses a product of
\emph{invariant link/color balls} and \(\chi_{\rm extra}=1\), or extra
indicators separately proved covariant. It is not asserted for arbitrary
scalar-coordinate boxes or arbitrary extra cutoffs. The exact real-window
proof precedes continuation. For a small real coarse gauge \(k\), physical
branch uniqueness and the block-constant lift give
\(V_*(W^k)=V_*(W)^{k^{\rm up}}\). With \(R_k\) the source-endpoint
adjoint on relative fields and \(D_k\) its retained restriction,
constraint equivariance and chart uniqueness imply
\[
 \Phi_{V_*(W^k)}(D_kz)=R_k\Phi_{V_*(W)}(z).
\]
Each retained adjoint is orthogonal and has determinant one. Hence it
preserves the sharp product of color balls exactly. Input gauge
invariance, conjugation invariance of the absolute contour penalty and
the actual Haar/pivot change of variables prove equality of the
\emph{literal real retained integrals}. No cutoff boundary term is
omitted. These are the exact identities proved in the preceding local
macrostep covariance argument; the analytic density then extends the
differentiated Ward identity over its connected local ball. Constant
conjugation follows by the same argument. The density version needs no
new logarithm branch. For other cutoffs this application is conditional,
while Theorem~\ref{fps-positive-split} remains unconditional.

Choose \(\theta\) supported in \(D_\epsilon\),
positive in its half ball. If \(W\) is in the half ball, then
\(m_\theta(W)>0\), and every contributing \(W^v\) lies in
\(D_\epsilon\). Consequently
\begin{equation}
 a_\mu(W)=\ell_\mu(W)\qquad\hbox{near the identity}.
 \label{fps-germ-equality}
\end{equation}
Thus a real measurable global good measure has exactly the previously
computed analytic germ. The graph-dependent shrink is explicit;
common-ball Ward continuation alone is not invariance between arbitrary
orbit endpoints whose connecting paths may leave that ball.

\subsection{Full histories, sources and raw normalization}
\label{fps-histories-section}

At each finite stage fix the map, gauge penalty, positive coupling,
window and output cutoff. Write
\(\mathcal T_j^{\rm good}=\mathcal A_j\) and
\(\mathcal T_j^{\rm rem}=\mathcal R_j\).
Starting with the original unnormalized finite Wilson measure
\(\mu_0=e^{-\gamma_0S_W}m_E\), define for every binary word \(h\)
\[
 \mu_h=\mathcal T_{n-1}^{h_{n-1}}\cdots
                         \mathcal T_0^{h_0}\mu_0.
\]
Every history is a positive invariant finite measure. Repeatedly using
Theorem~\ref{fps-positive-split} proves
\begin{equation}
 (F_{n-1}\circ\cdots\circ F_0)_*\mu_0
     =\sum_{h\in\{{\rm good},{\rm rem}\}^n}\mu_h,
 \qquad Z_W=\sum_h\mu_h(1).
 \label{fps-history-identity}
\end{equation}
The same equality holds against every bounded terminal coarse
observable. It is exact and includes all mixed histories, singular
parts and raw vacuum factors; it supplies no smallness estimate for
the remainder.

For bounded gauge-invariant original observables, the same linear
identities apply to their invariant signed numerator measures, by
positive/negative decomposition. Equivalently apply them to
\(e^{\sum t_iO_i}\mu_0\) on real compact source intervals and
differentiate, holding step maps and windows fixed. This provides
exact finite numerator and denominator bookkeeping, not uniform
observable estimates in a refinement limit.

To identify an all-good germ with the previously established
homogeneous map, its input density must actually equal
\(Z_j e^{-\gamma_jS_{M_j}+E_{j,\rm raw}}\), with \(\gamma_j=g_j^{-2}\), on a real neighborhood
\(D_{\rm in}\), and one must verify the literal containment
\begin{equation}
 e^{i\Phi_{V_*(W)}(g_j\xi)}V_*(W)\in D_{\rm in}
 \quad\hbox{for every retained \(\xi\) and the chosen \(W\)}.
 \label{fps-literal-containment}
\end{equation}
Every additional indicator is retained, unless it is proved identically
one there. Under this condition \eqref{fps-gaussian-density} is the
actual homogeneous local map, and \eqref{fps-germ-equality} embeds that
germ as a positive full-measure submeasure. For a fixed finite number
of stages, sufficiently weak couplings and small angle windows can meet
these finite margins: a logarithmic threshold is compatible with
\(gT(g)\to0\). This is not a volume-uniform assertion about the tuned
histories constructed earlier.

With \(A_0\) the flat retained precision, the exact field-independent
factor extracted in one stage is
\begin{equation}
 N_j=c(g_j)^{-n_s}h_0^{n_e-n_c}g_j^{n_z}(2\pi)^{n_z/2}
                       D_R(1)\det(A_0)^{-1/2}.
 \label{fps-vacuum-factor}
\end{equation}
The local output is
\(Z_jN_j e^{-\gamma_jS_{M_{j+1}}+F_{j,\rm raw}}\).
The raw correction contains the inherited background, actual pivot and
Gaussian determinant ratios and the exact retained expectation; its
value at zero is not discarded. Since \(S_{M_{j+1}}(0)=0\), if
\[
 F_{j,\rm raw}=\beta_jS_{M_{j+1}}+E_{j+1,\rm raw}+c_j,
 \qquad c_j=F_{j,\rm raw}(0)-E_{j+1,\rm raw}(0),
\]
then the exact bookkeeping is
\begin{equation}
 \gamma_{j+1}=\gamma_j-\beta_j,\qquad
 Z_{j+1}=Z_jN_j e^{c_j}.
 \label{fps-exact-extraction}
\end{equation}
Replacing an input representative by \(E'_{\rm raw}=E_{\rm raw}+c\)
requires \(Z'=Ze^{-c}\). Linearity of the orbit extension preserves
all these scalars in the globally defined good submeasure. No global
analytic classical action is required for these local identities.

\subsection{The quantitative interface remains separate}
\label{fps-boundary-section}

Let \(\nu\) be the full terminal pushforward and \(\nu_{\rm good}\)
its all-good submeasure. If the latter has positive mass, put
\(\epsilon_{\rm bad}=1-\nu_{\rm good}(1)/\nu(1)\).
Positivity and the mixture identity give
\begin{equation}
 \left|\frac{\nu(f)}{\nu(1)}-
             \frac{\nu_{\rm good}(f)}{\nu_{\rm good}(1)}\right|
 \le2\|f\|_\infty\epsilon_{\rm bad}.
 \label{fps-mixture-bound}
\end{equation}
No useful bound on \(\epsilon_{\rm bad}\) has been proved here. In large
volume, global all-good mass may be a poor criterion; a spatial
mixed-history expansion can be necessary for local observables.

The remaining tasks include uniform literal chart/orbit-domain matching,
interacting coercivity and normalization estimates for the actual
remainders, collared mixed-field factorization with the correct
conditional means and compressed precisions, and observable bounds with
the same full denominator. Homogeneous Gaussian bad-block bounds do not
alone prove them. The present result is an exact positive finite bridge
and an explicit remainder/history operator, not a completed global RG,
continuum measure, reconstruction theorem or physical mass gap.


% No preamble. Prefix fwp-. Full positive Wilson measure, 17 October 2026.
\section{A full-measure Wilson Peierls bound}
\label{fwp-section}

We prove a rare-event estimate under the actual normalized positive
Wilson measure. It is independent of the Gaussian small-field analysis.
The event controlled here is large plaquette curvature; identifying it
with the complement of the actual retained-coordinate window is a
separate, unresolved interface.

Let \(G=SU(N)\), \(m_G=\dim G\), \(\tau=\operatorname{Tr}/N\), and
\[
 \Lambda=\prod_{i=1}^d\mathbb Z/n_i\mathbb Z,\qquad
 n_i=2m_i b,\quad m_i\ge1,\quad b\ge3.
\]
There are \(V=|\Lambda|\) sites, \(dV\) independent links and
\(\binom d2V\) elementary unoriented plaquettes. Normalized Haar measure
is used on every link. Set
\begin{equation}
 h(W)=1-\operatorname{Re}\tau(W)
       =\tfrac12\|W-I\|_\tau^2,\qquad
 A(U)=\sum_p h(U_p),\qquad
 \mu_g(dU)=Z_g^{-1}e^{-A(U)/g^2}\prod_e dU_e,
 \label{fwp-measure}
\end{equation}
where \(\|X\|_\tau^2=\tau(X^*X)\).
All side lengths are at least six, avoiding short-period graph
degeneracies.

Fix any translated grid of cubes of side \(b\).
Let \({\cal P}_j^{\rm int}\) be the plaquettes whose vertices are strictly
inside cell \(j\), and define
\begin{equation}
 {\cal B}_j(\delta)
 =\left\{\max_{p\in{\cal P}_j^{\rm int}}h(U_p)\ge\delta^2\right\}.
 \label{fwp-event}
\end{equation}
Strict interior means that every coordinate of every used link endpoint
lies strictly between the corresponding walls.

\begin{theorem}[Prescribed bad cells]
\label{fwp-main}
There are \(g_0>0\) and \(C_{G,d}<\infty\), depending only on the group
and dimension, such that for \(0<g\le g_0\) and every set \(S\) of
distinct cells,
\begin{equation}
 \mu_g\left(\bigcap_{j\in S}{\cal B}_j(\delta)\right)
 \le z_b(g,\delta)^{|S|},
 \quad
 z_b(g,\delta)=
 \min\left\{1,
 \exp\left[-\frac{\delta^2}{g^2}
     +b^d\{d m_G\log(1/g)+C_{G,d}\}\right]\right\}.
 \label{fwp-main-bound}
\end{equation}
Constants are uniform in volume and in the selected cells. The source
regulators \(n=2L^M\), \(L\ge3\) odd, admit the construction with
\(b=L^u\) for each fixed \(u\ge1\) and \(M\ge u\).
\end{theorem}

\subsection{Exact site-plane slabs and the normalized chessboard inequality}
\label{fwp-slab-section}

Every cell wall is an integer coordinate, or \emph{site}, plane.
Its reflection preserves the lattice. An elementary plaquette cannot
have vertices on both strict sides of such a plane. It either lies in
the plane or belongs wholly to one adjacent closed half-space, possibly
touching the plane. Conditional on tangential links in the two opposite
reflection planes of the torus, the two open half configurations have
reflected product Haar measures and reflected Wilson weights.
The reflection quadratic form is therefore an integral of an absolute
square against a nonnegative boundary weight. This proves the needed
site-plane reflection positivity directly.

To obtain the chessboard inequality for an arbitrary even number of
cells, cut first in one coordinate into \(p=2m\) slabs of width \(b\).
Let \({\cal H}\) be \(L^2\) of the product Haar measure on tangential
links in one boundary plane. Integer translation identifies successive
boundary spaces.

For a bounded slab observable \(f\), supported strictly inside the
slab in the cut direction, form the integral kernel \(K_f\) by integrating
all nonboundary slab links. Give it the full action of every plaquette
assigned to the slab, including plaquettes touching a boundary but not
lying in it, and half the action of each plaquette lying entirely in
either boundary. Multiply by \(f\).
The compact integration domain and bounded integrand make \(K_f\)
Hilbert--Schmidt.

Gluing the kernels by boundary Haar integration gives exactly
\begin{equation}
 \operatorname{Tr}(K_{f_0}\cdots K_{f_{p-1}})
       =Z_g\!\left[\prod_{j=0}^{p-1}f_j\right].
 \label{fwp-gluing}
\end{equation}
Every boundary link is integrated once; every boundary plaquette gets
two halves; every other link and plaquette is assigned once.
No gauge-volume factor occurs. Reflection exchanges the two boundaries,
preserves \(h(U_p)\), and inverts a chosen link coordinate when its
orientation reverses. Haar invariance gives
\begin{equation}
 K_f^*=K_{\theta\overline f}.
 \label{fwp-adjoint}
\end{equation}
For real indicators complex conjugation has no effect.

Schatten H\"older, with the even integer \(p\), gives
\begin{equation}
 \left|\operatorname{Tr}\prod_{j=0}^{p-1}K_{f_j}\right|
 \le\prod_{j=0}^{p-1}
 \left[\operatorname{Tr}\bigl((K_{f_j}^*K_{f_j})^{p/2}\bigr)
                                                  \right]^{1/p}.
 \label{fwp-holder}
\end{equation}
Each factor is the exact Wilson partition integral with \(p\) alternating
reflected copies of one slab. The torus period remains \(pb\).
This explains precisely why the even cell count is required.

Apply \eqref{fwp-holder} successively in all coordinate directions.
A slab source is the product of its constituent cell sources.
Orthogonal affine coordinate reflections commute on the unoriented
graph, with the corresponding orientation inversions on link variables.
The final factor associated with an original cell is therefore its
full reflected dissemination. Intersections of cut planes require no
extra normalization: every step is the exact one-axis gluing
\eqref{fwp-gluing}.

Write \(N_c=V/b^d\) for the number of cells. Let \(\Theta_j\) map the
reference cell to cell \(j\) by the alternating reflection pattern;
even counts make this pattern periodic. For a reference-cell observable
\(f\), write
\({\cal D}(f)=\prod_i\Theta_i f\).
Pull each original observable back to this reference cell.
The result of the iteration is
\begin{equation}
 \left|Z_g\left[\prod_j\Theta_j f_j\right]\right|
 \le\prod_j Z_g[{\cal D}(f_j)]^{1/N_c}.
 \label{fwp-chessboard}
\end{equation}
For nonnegative indicators, all terms are nonnegative. Dividing by the
\emph{full} partition function, and using \(f_j=1\) off a selected set,
gives
\begin{equation}
 \mu_g\left(\prod_{j\in S}\Theta_j f_j\right)
 \le\prod_{j\in S}
       \left(\frac{Z_g[{\cal D}(f_j)]}{Z_g}\right)^{1/N_c}.
 \label{fwp-normalized}
\end{equation}
In particular, the normalization is not an independently chosen
small-field partition function.

The geometry used here is simpler than a diagonal fold: all walls are
site planes and no alternating site/link scheme is needed.
For \(n=2L^M\), \(b=L^u\), the number of cells in every direction is
\(2L^{M-u}\), which is even despite \(L\) being odd.

\subsection{A lower bound for the full normalization}
\label{fwp-normalization-section}

There are \(\epsilon_0>0\) and \(c_G\in(0,1]\) such that
\begin{equation}
 \operatorname{Haar}\{U:\|U-I\|_\tau\le\epsilon\}
       \ge c_G\epsilon^{m_G},\qquad 0<\epsilon\le\epsilon_0.
 \label{fwp-haar-ball}
\end{equation}
Indeed the Haar density in a fixed local exponential chart is smooth
and strictly positive. This is a finite-dimensional group estimate.

If all links lie in that \(\epsilon\)-ball, a four-term product
telescoping and unitarity show
\[
 \|U_p-I\|_\tau\le4\epsilon,\qquad h(U_p)\le8\epsilon^2.
\]
For \(g\le g_0:=\min(\epsilon_0,1)\), restrict the full integral to
\(\|U_e-I\|_\tau\le g\) on every link. It follows that
\begin{equation}
 Z_g\ge(c_Gg^{m_G})^{dV}
             \exp\!\left[-8\binom d2V\right],
 \qquad
 -\frac1V\log Z_g
       \le d m_G\log(1/g)+C_{G,d},
 \label{fwp-z-lower}
\end{equation}
where one can take
\[
 C_{G,d}=8\binom d2+d\log(1/c_G).
\]
This is only a lower bound for the normalization. It does not claim
that the small-link box has large probability.

On the full dissemination of \({\cal B}_j(\delta)\), every cell
contains a distinct interior plaquette of energy at least \(\delta^2\).
Thus \(A(U)\ge N_c\delta^2\) there and
\begin{equation}
 Z_g[{\cal D}(\mathbf1_{\cal B})]
       \le e^{-N_c\delta^2/g^2}.
 \label{fwp-disseminated-energy}
\end{equation}
The remaining product Haar integral is at most one.
Combining \eqref{fwp-z-lower} and \eqref{fwp-disseminated-energy}
bounds each factor in \eqref{fwp-normalized} by the exponential in
\eqref{fwp-main-bound}; it is also at most one.
This proves Theorem~\ref{fwp-main}.

More generally, a cell event forcing its strictly interior action to be
at least \(e_j\) has the factor
\begin{equation}
 \min\left\{1,\exp\left[-e_j/g^2
              +b^d\{d m_G\log(1/g)+C_{G,d}\}\right]\right\}.
 \label{fwp-energy-general}
\end{equation}
A threshold \(\|U_p-I\|_\tau\ge\delta\) uses \(e_j=\delta^2/2\).
At fixed \(b\), the threshold \(\delta(g)=gR(g)\) is rare whenever
\(R(g)^2/\log(1/g)\to\infty\), for example
\(R(g)=(\log(1/g))^a\) with \(a>1/2\).
The explicit \(b^d\) cost precludes claiming the same constants when
the cell size grows with the cutoff.

\subsection{Prescribed plaquettes and Peierls clusters}
\label{fwp-color-section}

Let \(S\) be any prescribed set of distinct elementary plaquettes.
Color \((x;\mu,\nu)\), \(\mu<\nu\), by \((x\bmod b,\mu,\nu)\).
There are
\begin{equation}
 Q_b=b^d\binom d2
 \label{fwp-color-count}
\end{equation}
colors. Choose deterministically a color occurring at least
\(|S|/Q_b\) times, with anchor residue \(a\), and shift the cell grid
by \(a-(1,\ldots,1)\). Every selected anchor has relative coordinates
\((1,\ldots,1)\); its plaquette vertices have coordinates \(1\) or \(2\)
and hence are strictly interior for \(b\ge3\).
Distinct anchors of the fixed color lie in distinct cells.
The fixed orientation prevents two selected plaquettes at one anchor.
The all-bad event on \(S\) implies the bad-cell event on this subset.
Therefore
\begin{equation}
 \mu_g\{h(U_p)\ge\delta^2\text{ for all }p\in S\}
       \le z_b(g,\delta)^{|S|/Q_b}.
 \label{fwp-arbitrary-plaquettes}
\end{equation}
This is a finite deterministic coloring argument, not an independence
assumption. The chessboard proof holds for each of the translated grids.

For the nearest-neighbor cell graph, of degree at most \(2d\), the
number of connected sets of size \(m\) containing a fixed cell is at
most \((2d)^{2m}\): choose a canonical spanning tree and its depth-first
traversal. The traversal has length \(2(m-1)\) and determines the visited
set. A cluster of size at least \(m\) contains such a connected subset
of size exactly \(m\). Consequently
\begin{equation}
 \mu_g\{|\text{bad cluster at }j|\ge m\}
       \le[(2d)^2z_b(g,\delta)]^m.
 \label{fwp-cell-cluster}
\end{equation}
For any fixed finite-degree plaquette adjacency graph the same argument
and \eqref{fwp-arbitrary-plaquettes} apply. For adjacency by shared
vertices, \(\Delta_d=16\binom d2\) is a safe degree bound, giving
\begin{equation}
 \mu_g\{|\text{bad plaquette cluster at }p|\ge m\}
       \le[\Delta_d^2z_b(g,\delta)^{1/Q_b}]^m.
 \label{fwp-plaquette-cluster}
\end{equation}
Periodic identifications do not increase these counting bounds.
These are intersection and cluster estimates, not a claim of full
stochastic domination by a product measure.

The union estimate for any bad cell is only \(N_c z_b\).
Thus the theorem does not assert that an arbitrarily large torus lies
entirely in one small-field sector with uniformly high probability.

\subsection{The remaining large-field interface}
\label{fwp-scope}

The distinction between curvature and retained-coordinate windows is
essential. For a link with distinct endpoints, gauge invariance at one
endpoint makes its marginal left-invariant under \(G\). Hence
\begin{equation}
 \mu_g\{U_e\in A\}=\operatorname{Haar}(A)
 \label{fwp-haar-marginal}
\end{equation}
for every measurable \(A\), independently of \(g\).
An unfixed link is therefore outside a shrinking identity window with
probability tending to one. Pure-gauge fields can have all plaquettes
equal to \(I\) and individual links far from \(I\).

The actual retained fluctuation \(\xi\) instead depends on the nonlinear
constraint, its minimizing background, the contour gauge and the pivot
chart. No result above proves that \(\|\xi\|>R\) forces the interior
Wilson action to exceed \(c g^2R^2\). Such a deterministic, quantitative
implication for the exact chart would be an additional input.
A local tree-gauge estimate from small plaquettes does not by itself
identify its potential with \(\xi\), or establish the required
\(C^{1,\alpha}\) background domain.

Global holonomy is a second obstruction. For example,
\(U_1(x)=e^{i\theta H/n_1}\), \(U_\mu(x)=I\) for \(\mu\ne1\),
has every plaquette equal to \(I\), while its noncontractible Wilson
line is \(e^{i\theta H}\). Plaquette rarity cannot simply discard these
sectors.

Finally, the proof uses the positive bare Wilson measure. It has not
been extended here to complex source measures, arbitrary conditioned
boundaries or effective actions after many RG steps. It supplies a
genuine full-measure large-field input: uniform prescribed-event and
cluster bounds on the original regulator. It does not yet supply the
exact retained-window implication, a convergent sum over all large-field
components with analytic source and locality bounds, or a complete
Wilson RG history.

For context, the general reflection/chessboard framework is
Fr\"ohlich--Israel--Lieb--Simon, \emph{Communications in Mathematical
Physics} \textbf{62} (1978), 1--34,
DOI \(10.1007/BF01940327\); gauge positivity is treated by
Osterwalder--Seiler, \emph{Annals of Physics} \textbf{110} (1978),
440--471, DOI \(10.1016/0003\text{-}4916(78)90039\text{-}8\).
The site-plane gluing and normalization used here were proved explicitly
above rather than imported with an unspecified boundary convention.



\section{A normalized conditional Wilson tail in a fixed box}
\label{fbw:section}

This section controls the entire compact Wilson integral for a fixed
box with small boundary data. The constants depend on the box. It
does not identify this box gauge with the nonlinear constrained
coordinates of the homogeneous renormalization step.

Fix a cubical lattice box \(\Omega=\{0,\ldots,R\}^d\), \(d=4\),
\(R\geq2\), with its full cubical cell structure. Let \(E_\partial\)
be the edges in its topological boundary and \(E_\circ\) the remaining
edges. The boundary edge graph is connected. Fix
\(G=SU(N)\), with faithful normalized trace \(\tau=\operatorname{Tr}/N\),
and put
\[
 h(U)=1-\operatorname{Re}\tau(U)
     =\tfrac12\norm{U-I}_\tau^2,\qquad
 S_V(U)=\sum_{p\subset\Omega}h(U_p).
\]
Boundary edge variables \(V\in G^{E_\partial}\) are prescribed.
All remaining edges have normalized Haar measure. Gauge transformations
are fixed to the identity at every boundary vertex.

Choose a spanning forest \(\mathcal F\) of the full vertex set with
exactly one boundary vertex in every component, and all interior
vertices included. Its edges lie in \(E_\circ\). Such a forest exists:
start with a spanning tree of the connected boundary graph, extend
it to a spanning tree of the whole graph, and remove its boundary
tree edges. There is a unique allowed gauge transformation making
every forest edge the identity. Write \(y\in\mathcal Y=G^{E_\circ
\setminus\mathcal F}\) for the remaining edge variables in that gauge.
The number of real coordinates is
\[
 n=(\dim G)\,|E_\circ\setminus\mathcal F|.
\]
For this box \(n>0\). Fix an orthonormal real Lie algebra basis and
write \(y=\exp(ix)\) componentwise on a neighborhood of the identity.
Here \(\norm{x}_\infty=\max_e\norm{x_e}_\tau\) uses the
adjoint-invariant Lie algebra norm on each edge.

\begin{theorem}[Full conditional tail at fixed geometry]
\label{fbw:theorem}
There are a neighborhood \(\mathcal V\) of the identity boundary
configuration, constants \(g_0,r,c,C>0\), and a real analytic
function \(x_*(V)\) on \(\mathcal V\), with \(x_*(I)=0\), such that
the following holds. The forest-gauge action has the unique global
minimum \(y_*(V)=\exp(ix_*(V))\). For
\(0<g\leq g_0\) and \(1\leq T\leq r/g\), its normalized Wilson
measure satisfies
\begin{equation}
 \frac{\displaystyle
    \int_{\mathcal Y\setminus W_{V,g,T}}
           e^{-S_V(y)/g^2}\,dy}
      {\displaystyle\int_{\mathcal Y}e^{-S_V(y)/g^2}\,dy}
 \leq C e^{-cT^2},
 \qquad
 W_{V,g,T}=\{\exp(ix):\norm{x-x_*(V)}_\infty\leq gT\}.
 \label{fbw:tail}
\end{equation}
The displayed window is contained in a single fixed exponential
coordinate neighborhood. The same ratio is the probability of the
corresponding gauge-fixed event under the original, ungauged box
Wilson measure.
\end{theorem}

\begin{proof}
We give the global normalization argument, including the gauge and
topological points.

\emph{Exact forest gauge.}
Orient each forest component away from its boundary root. Recursively
solving the transformation equation along its edges gives the unique
interior gauge values. Changing from the original forest edge variables
to these gauge values is triangular; each step is a left or right
translation, or inversion, on \(G\), hence preserves normalized Haar
measure. For fixed gauge values the other edge transformations preserve
Haar independently. The change of variables is therefore an exact
product-Haar identity. A gauge-invariant action and an event defined
through the unique forest representative are independent of the
interior gauge variables. Their integrals over those variables equal
one. No determinant or volume factor is discarded. Boundary edges
are unchanged even when \(V\ne I\); this is why a forest rooted at
all boundary vertices is used.

\emph{The unique zero and its Hessian.}
At \(V=I\), zero action means every plaquette holonomy is \(I\).
The cubical box is simply connected, so a flat edge field is a pure
gauge: parallel transport from a fixed vertex is independent of the
path, since elementary plaquette moves generate all path homotopies.
Identity boundary edges force the vertex gauge to be constant on the
connected boundary; normalize that constant to \(I\). The field
thus has an allowed gauge transformation to the identity. Uniqueness
of the forest representative gives \(y=I\). Consequently the
continuous nonnegative function \(S_I\) has exactly one zero on the
compact space \(\mathcal Y\).

For a Lie algebra edge field \(a\), zero on boundary and forest edges,
the quadratic part is
\[
 \tfrac12\sum_{p\subset\Omega}\norm{(da)_p}_\tau^2.
\]
If it vanishes, the elementary cellular Poincare lemma on a box gives
\(a=d\phi\). Boundary connectedness lets us choose \(\phi=0\) at
every boundary vertex. Vanishing on each forest edge then gives
\(\phi=0\) at every vertex. Hence \(a=0\). The Hessian in the \(n\)
remaining coordinates is positive definite. This is a finite
dimensional assertion; no uniform lower eigenvalue as \(R\) varies
is asserted.

\emph{Uniform local well and global gap.}
The analytic implicit-function theorem supplies a critical point
\(x_*(V)\) near zero for \(V\) close to \(I\). Choose a coordinate
ball about zero on which the Hessian at all such \(V\) is uniformly
positive. On a smaller fixed ball about \(x_*(V)\), Taylor's theorem
gives constants \(c_0,C_0>0\) with
\begin{equation}
 c_0|x-x_*(V)|^2\leq S_V(\exp(ix))-S_V(y_*(V))
       \leq C_0|x-x_*(V)|^2.
 \label{fbw:quadratic}
\end{equation}
The balls can be chosen uniformly after shrinking \(\mathcal V\).
To see a global gap, first choose a still smaller ball about zero.
On the compact complement its \(V=I\) action has a strictly positive
minimum. Uniform continuity of \(S_V\), and \(S_V(y_*(V))\to0\),
preserve half this gap for sufficiently small boundary data.
Together with strict convexity on the coordinate ball this proves
uniqueness of the global minimum and yields a fixed chart region
\(\mathcal U\subset\mathcal Y\) and \(\eta>0\) for which
\begin{equation}
 S_V(y)\geq S_V(y_*(V))+\eta\quad(y\notin\mathcal U).
 \label{fbw:gap}
\end{equation}
Choose \(\mathcal U\) small enough that the lower quadratic bound in
\eqref{fbw:quadratic} holds on all of it, and large enough that a
fixed ball about every \(x_*(V)\) lies in it. These choices are
compatible by first making the convexity ball larger than the
compact-complement ball and then shrinking \(\mathcal V\).
The product Haar coordinate density on its closure is bounded above
and below by positive constants.

\emph{Denominator and local numerator.}
Writing \(S_{\min}=S_V(y_*(V))\), integrate over the coordinate
ball \(|x-x_*|\leq g\), reducing \(g_0\) as needed:
\begin{equation}
 Z_V(g)\geq c_1 g^n e^{-S_{\min}/g^2}.
 \label{fbw:denominator}
\end{equation}
Within \(\mathcal U\), the complement of the window forces
\(|x-x_*|\geq gT\). The lower quadratic bound and the upper Haar
density bound give, for \(T\geq1\),
\begin{align}
 \int_{\mathcal U\setminus W_{V,g,T}}e^{-S_V/g^2}\,dy
 &\leq C_1 e^{-S_{\min}/g^2}
       \int_{|z|\geq gT}e^{-c_0|z|^2/g^2}\,dz \notag\\
 &\leq C_2 g^n e^{-S_{\min}/g^2}e^{-c_2T^2}.
 \label{fbw:localnumerator}
\end{align}
The elementary Gaussian tail estimate absorbs its finite dimensional
polynomial factor by reducing \(c_2>0\).

\emph{The rest of the compact group integral.}
The remaining numerator is at most
\(e^{-(S_{\min}+\eta)/g^2}\), since product Haar mass is one.
After \eqref{fbw:denominator} its contribution is at most
\(C_3g^{-n}e^{-\eta/g^2}\). Since
\[
 \sup_{0<g\leq g_0}g^{-n}e^{-\eta/(2g^2)}<\infty,
\]
and \(T\leq r/g\), choosing \(r\) small and
\(c\leq \min\{c_2,\eta/(2r^2)\}\) bounds this contribution by
\(C_4e^{-cT^2}\). Reduce \(r\) also to keep the entire sup-norm
window inside the coordinate chart. Combining the two numerators
with the denominator proves \eqref{fbw:tail}.
\end{proof}

\begin{corollary}[Logarithmic windows]
\label{fbw:logwindow}
For fixed \(a>1/2\), take \(T(g)=(\log(1/g))^a\). For every
\(m>0\) there are \(g_m,C_m>0\) such that
\[
 \mu_{V,g}(\mathcal Y\setminus W_{V,g,T(g)})\leq C_m g^m
 \quad(0<g\leq g_m,\ V\in\mathcal V).
\]
\end{corollary}
\begin{proof}
Here \(gT(g)\to0\) and \(T(g)^2/\log(1/g)\to\infty\).
Apply \eqref{fbw:tail}.
\end{proof}

\paragraph{Conditional interpretation.}
Embed the box in a larger lattice without local periodic
identifications. A plaquette containing an edge of \(E_\circ\)
lies entirely inside the box: an edge not in the boundary has all
its transverse coordinates strictly inside the corresponding
intervals. Thus, after the exterior and boundary edges are fixed,
the part of the bare Wilson action depending on the integrated edges
is exactly \(S_V\), up to an additive constant. Such constants cancel
in the normalized law. The theorem is therefore a statement about
the actual conditional bare Wilson measure whenever its boundary
field belongs to \(\mathcal V\). A prescribed boundary gauge
transformation bringing the data into \(\mathcal V\) also permits
the theorem after transforming the event consistently.

\paragraph{The precise interface left open.}
This result controls a fixed-box forest-coordinate window about
the unconstrained conditional minimizer. It does not control an
arbitrary boundary field, a box growing with the cutoff, a
nonlinear coarse-average fiber, or a mixed-history effective action
whose global stability has not been proved. In particular an
unconditional probability bound cannot simply be conditioned on a
coarse-map value with no lower bound for the corresponding fiber
normalization. The powers \(g^n\) cancel here because both integrals
are estimated in the same compact conditional model. Extending
that cancellation to the actual constrained local map is a
remaining analytic task.


\section{From curvature rarity to full-measure box windows}
\label{fbb:section}

The two preceding probability estimates can be joined without
assuming that a boundary field is typically close to the identity
in an unfixed gauge. The resulting windows are defined by an
explicit local tree gauge and a conditional Wilson minimizer.
They still differ from the nonlinear coarse-average windows in
the homogeneous renormalization construction.

\begin{lemma}[A boundary-only gauge at fixed geometry]
\label{fbb:boundarygauge}
For the fixed four-dimensional cubical box \(\Omega\), choose a
spanning tree \(\mathcal T_\partial\) of its boundary edge graph
and a root vertex. There is a finite constant \(L_R\) such that
every boundary field has a boundary gauge, determined only by its
boundary edges, with
\[
 V_e=I\quad(e\in\mathcal T_\partial),\qquad
 \max_{e\in E_\partial}\norm{V_e-I}_\tau
 \leq L_R\max_{p\subset\partial\Omega}\norm{U_p-I}_\tau.
\]
The estimate concerns boundary plaquettes, not plaquettes through
the interior. No assertion uniform in \(R\) is made.
\end{lemma}
\begin{proof}
Fix the root gauge to \(I\) and gauge the boundary tree edges to \(I\)
by recursion. Every remaining boundary edge becomes the holonomy
of its fixed tree fundamental loop. The boundary cubical complex is
homeomorphic to \(S^3\); its two-skeleton has trivial fundamental
group. Each of the finitely many fundamental loops is therefore a
finite product of conjugates of elementary boundary plaquette
loops and their inverses, as a word in the graph edge generators.
This assertion follows directly from the cellular presentation of
the fundamental group: the plaquette boundary words normally
generate all relations in the graph fundamental group. Choose one
such finite word once for each non-tree edge and let \(L_R\) be the
largest number of plaquette factors used.

For unitary matrices \(A_j\), telescoping gives
\(\norm{\prod_j A_j-I}_\tau\leq\sum_j\norm{A_j-I}_\tau\).
Conjugation and inversion preserve this distance. Applying this
inequality to the chosen word proves the bound, regardless of
the size of its conjugating holonomies.
\end{proof}

Choose \(\varepsilon_R>0\) so small that
\(\sqrt2 L_R\varepsilon_R\) lies inside the small boundary
neighborhood in Theorem~\ref{fbw:theorem}. We may and do take that
neighborhood to be a conjugation-invariant product ball. Call a
box boundary good when
\[
 h(U_p)<\varepsilon_R^2\quad\hbox{for every }p\subset\partial\Omega.
\]
On a good boundary, apply the boundary-only gauge in
Lemma~\ref{fbb:boundarygauge}; extend its vertex values into the
box arbitrarily, then perform the rooted forest gauge of
Section~\ref{fbw:section}, keeping the transformed boundary fixed.
The final representative is independent of that arbitrary extension:
two extensions differ by an interior gauge, which the unique forest
normalization removes. It is also the representative obtained by
gauging the full spanning tree
\(\mathcal T_\partial\cup\mathcal F\) to \(I\).

Define the box bad event \(D_\Omega(g,T)\) as follows:
the boundary is bad, or the boundary is good and the resulting
interior representative lies outside \(W_{V,g,T}\). All these maps
are measurable. The event is gauge invariant. Indeed, an arbitrary
original gauge transformation changes the full tree representative
only by constant conjugation at its root. The conditional minimizer
transforms by that same conjugation, by uniqueness, and the local
logarithm and edge-max norm respect conjugation.

\begin{theorem}[Prescribed local box defects in the full Wilson measure]
\label{fbb:boxes}
Fix \(R,b\), with \(b\geq3\), and consider a periodic four-dimensional
lattice with side lengths divisible by \(2b\). Let
\(\Omega_1,\ldots,\Omega_m\) be embedded translates of the fixed box
with pairwise disjoint vertex sets. Put
\[
 Q_b=b^4\binom42,\qquad
 N_\partial=\#\{p:p\subset\partial\Omega\}=24R^3,
\]
and let \(z_b(g,\varepsilon_R)\) be the bad-cell factor of
Section~\ref{fwp-section}. There are constants \(C,c,r,g_0>0\),
independent of the total torus size and of the chosen boxes, such
that
\begin{equation}
 \mu_g\left(\bigcap_{i=1}^m D_{\Omega_i}(g,T)\right)
 \leq \left[
       C e^{-cT^2}
       +N_\partial z_b(g,\varepsilon_R)^{1/Q_b}
       \right]^m
 \quad(0<g\leq g_0,\ 1\leq T\leq r/g).
 \label{fbb:joint}
\end{equation}
In particular, for every \(a>1/2\), the bracket with
\(T(g)=(\log(1/g))^a\) is smaller than every fixed power of \(g\)
as \(g\downarrow0\).
\end{theorem}
\begin{proof}
Condition on every edge outside the union of the boxes' interior
edge sets \(E_{\circ,i}\). This includes all box boundary edges.
Every plaquette meeting \(E_{\circ,i}\) lies wholly in \(\Omega_i\).
Disjointness of the vertex sets therefore ensures that no plaquette
meets two integrated edge sets. The conditional law factorizes
exactly into the independent conditional Wilson laws of the
individual boxes; all remaining action terms are constant under
these integrations.

Write \(B_i\) for a bad boundary and \(F_i\) for the interior-window
failure on a good boundary. Each \(B_i\) is measurable under the
conditioning. The boundary gauge is also fixed under it, and it
preserves the conditional product Haar law and the action.
Theorem~\ref{fbw:theorem} consequently gives
\[
 \mathbb P(F_i\mid\hbox{conditioned edges})
 \leq a_g\,\mathbf1_{B_i^c},\qquad a_g=Ce^{-cT^2}.
\]
The \(F_i\)'s are conditionally independent.

For any subset \(J\) of boxes, a union bound over one offending
boundary plaquette in each box, followed by the prescribed-plaquette
bound of Section~\ref{fwp-section}, gives
\begin{equation}
 \mu_g\left(\bigcap_{i\in J}B_i\right)
 \leq \left[N_\partial
              z_b(g,\varepsilon_R)^{1/Q_b}\right]^{|J|}.
 \label{fbb:boundaryrare}
\end{equation}
The selected plaquettes are distinct because the boxes have disjoint
vertex sets. This is a correlated full-measure estimate; independence
of the boundaries is not used.

Expand \(\bigcap_i(B_i\cup F_i)\) over \(J\), the boxes assigned
the boundary-bad alternative. Conditional factorization bounds
the contribution for \(J\) by
\(a_g^{m-|J|}\mu_g(\bigcap_{i\in J}B_i)\), after dropping other
good-boundary indicators. Substituting
\eqref{fbb:boundaryrare} and summing the binomial expression proves
\eqref{fbb:joint}.

Finally, the explicit bad-cell exponent has fixed threshold
\(\varepsilon_R\) and fixed \(b\):
\[
 z_b(g,\varepsilon_R)^{1/Q_b}
 \leq \exp\left[-\frac{\varepsilon_R^2}{Q_b g^2}
                       +C_{b,G}\log(1/g)+C_{b,G}\right]
\]
for sufficiently small \(g\). This and
Corollary~\ref{fbw:logwindow} prove the last assertion.
\end{proof}

\paragraph{What this joins, and what it does not.}
The theorem gives joint rarity of actual, gauge-invariant local
coordinate defects under the complete bare Wilson measure. It
combines a correlated boundary estimate with conditional interior
normalization and does not replace the Wilson measure by a Gaussian
law. For a fixed separated grid of such boxes, the usual connected-set
counting yields Peierls tails with the displayed bracket.

The representative and minimizer are those of a fixed conditional
box with all its boundary edges held fixed. They are not identified
with the averaged, constrained minimizer and independent retained
cutoffs of the existing multiscale map. No derivative regularity,
complex-source estimate, growing-box constant, or post-renormalization
reflection positivity follows from this theorem. Thus it is a
concrete full-measure local bridge, while the multiscale mixed-history
induction remains an additional theorem to prove.


% No preamble. Prefix fmg-. Primary interface and proved finite geometry.
\section{Local gauge coordinates and the complete-density interface}
\label{fmg-section}

The homogeneous local trajectory does not yet identify all sectors
of the full Wilson integral. This section records a precise published
interface, proves a finite local gauge lemma, and shows why that
lemma alone does not supply a uniformly smooth physical background.

\subsection{The strongest published complete-density input}
\label{fmg-source}

Ba{\l}aban's Theorem~1 in \emph{Commun.\ Math.\ Phys.}\ 122
(1989), p.~355, supplies the complete effective-density class of
\emph{Commun.\ Math.\ Phys.}\ 119 (1988), Section~2, provided the
effective couplings of that construction remain sufficiently small.
The density bounds are (0.1), pp.~355--356; its large-field
exponentiation concludes at (1.98)--(1.102), pp.~390--391.
The original articles are available at
\href{https://scispace.com/pdf/large-field-renormalization-ii-localization-exponentiation-h8b7ih93xo.pdf}
{CMP122-II} and
\href{https://scispace.com/pdf/convergent-renormalization-expansions-for-lattice-gauge-29jyexd38e.pdf}
{CMP119}.

The relevant CMP119 input locations are:
\[
\begin{array}{c|l}
\text{pages/equations}&\text{interface}\\ \hline
243\text{--}244,\ (0.1)\text{--}(0.2)&
       \text{full pushforward and the additional operation}\\
246\text{--}248,\ (1.1)\text{--}(1.10)&
       \text{partitions and collars}\\
254\text{--}257,\ (2.1)\text{--}(2.18)&
       \text{multiscale determining sets and histories}\\
258,\ (2.19)\text{--}(2.22)&
       \text{large-field operator factorization}\\
258\text{--}262,\ (2.23)\text{--}(2.42)&
       \text{action, analytic domains and boundary terms}.
\end{array}
\]
In particular (2.38), p.~261, requires a local gauge with potential
and first-derivative bounds, alongside the other multiscale
regularity/current conditions. It is not a conclusion from one
bare plaquette threshold. The theorem concerns the specified
complete-density construction, not every small analytic local map.
OCR-damaged numerical constants in these source formulas are not
retranscribed here.

To use that result for our exact \(F_R=M_\ell^s\) local step,
one must either verify its complete induction class for this
operator or prove a replacement large-field theorem. In particular:
the full compact-field pushforward and partition must be specified;
the good-sector term must equal our actual retained-window
integral; exterior and transition data must be retained; and
normalized complementary-region operations must be bounded.
The source's own coupling recursion must be identified or replaced
by an independently justified one. Uniform weakness of our tuned
numbers alone is not equality of the two recursions.

Our strict product-ball families have an invariant total.
They do not automatically supply individually invariant local
orbit/current extensions or uniform analytic estimates on
multiscale determining sets. Nor does composition of averaging
constraints identify a one-shot window with sequential primitive
windows. These are concrete transfer obligations, without an
assertion that the source theorem cannot be adapted.

\subsection{An explicit finite cubical tree gauge}
\label{fmg-tree}

Let \(G\) be a compact matrix Lie subgroup of \(U(r)\), with
operator norm. Let
\(\Lambda_n=\{0,\ldots,n\}^d\), \(d\ge2\), carry its full cubical
cell structure. Oppositely oriented bonds carry inverse matrices.
For every vertex \(x\), let \(p(x)\) be the path from zero which
first increases coordinate one, then coordinate two, and so on.
These paths form a rooted spanning tree.

\begin{lemma}[Finite cubical gauge bound]
\label{fmg-tree-lemma}
Suppose every elementary plaquette obeys
\begin{equation}
                  \|U(\partial p)-I\|\le\epsilon.
 \label{fmg-plaquette-bound}
\end{equation}
Set
\begin{equation}
 v(x)=U(p(x)),\qquad
 U^v(x,y)=v(x)U(x,y)v(y)^{-1}.
 \label{fmg-gauge}
\end{equation}
Every tree edge of \(U^v\) is the identity, and each positive
\(\mu\)-edge satisfies
\begin{equation}
 \|U^v(x,x+e_\mu)-I\|
       \le\left(\sum_{\nu>\mu}x_\nu\right)\epsilon
       \le(d-1)n\epsilon.
 \label{fmg-edge-bound}
\end{equation}
Choose a fixed local exponential-chart radius
\(0<\delta_G\le1\). If \((d-1)n\epsilon<\delta_G\), write
\(U^v_b=\exp(i\eta A_b)\) in its unique local logarithm chart.
For physical mesh \(\eta\) and side \(\ell=n\eta\),
\begin{equation}
 \|A\|_\infty
       \le 2(d-1)n\epsilon/\eta
       =2(d-1)\ell\epsilon/\eta^2 .
 \label{fmg-physical-bound}
\end{equation}
\end{lemma}

\begin{proof}
The transformed edge is the product around the loop consisting
of \(p(x)\), the additional \(\mu\)-edge and reverse \(p(x+e_\mu)\).
Move the additional edge backwards past the later coordinate
segments. There are exactly \(\sum_{\nu>\mu}x_\nu\) elementary
square exchanges. Each exchange multiplies the loop holonomy
by a conjugate of a plaquette or its inverse, and the final
boundary word cancels. Unitary conjugation and inversion
preserve distance to \(I\). Moreover
\[
       \|W_1\cdots W_m-I\|
                     \le\sum_{j=1}^m\|W_j-I\|
\]
for unitary matrices, by telescoping. This proves
\eqref{fmg-edge-bound}; tree edges require no exchanges.
On the chosen local chart,
\(\|\log W\|\le2\|W-I\|\), as follows from the scalar eigenangle
bound. Dividing by \(\eta\) proves \eqref{fmg-physical-bound}.
\end{proof}

Before the logarithm, this gauge is a finite product of link
matrices and inverses. With a strict chart margin it is real
analytic and has a holomorphic continuation near each admitted
compact configuration. No refinement-uniform complex radius
or derivative bound follows just from
\eqref{fmg-edge-bound}. On a general simply connected cell complex
the same argument costs the filling areas of the tree fundamental
loops; simple connectivity alone does not bound those areas.

For a fixed finite current box the lemma gives local logarithmic
coordinates from sufficiently small plaquettes. For a fixed
physical box under refinement, its sufficient potential threshold
is \(\epsilon=O(\eta^2 b/\ell)\). It does not yet identify relative
fluctuations about a prescribed constrained minimum.

\subsection{A gauge-invariant obstruction to physical smoothness}
\label{fmg-holder-obstruction}

Fix \(0<\alpha<1\), a small \(f>0\), and a nonzero Cartan
generator \(H\) with \(\|H\|=1\). On a fixed physical cube
with decreasing mesh \(\eta\), let \(x\) denote integer lattice
coordinates and put
\begin{equation}
 U_1(x)=\exp\!\bigl(i\eta^2 f x_2\chi(x_1)H\bigr),\qquad
 U_\mu(x)=I\quad(\mu\ne1),
 \qquad
 \chi(k)=\begin{cases}0&k\text{ even},\\1&k\text{ odd}.
                     \end{cases}
 \label{fmg-oscillating-field}
\end{equation}
All non-\(12\) plaquettes are \(I\), and
\begin{equation}
               U(\partial p_{12}(x))
                    =\exp(-i\eta^2 f\chi(x_1)H).
 \label{fmg-oscillating-curvature}
\end{equation}
Thus every plaquette is within \(\eta^2f\) of \(I\).

\begin{proposition}[No uniform \(C^{1,\alpha}\) gauge]
\label{fmg-no-smooth-gauge}
For every fixed \(B<\infty\), sufficiently fine configurations
\eqref{fmg-oscillating-field} admit no gauge representative
\(U^v=\exp(i\eta A)\) on the cube satisfying
\begin{equation}
 \|A\|_\infty+\|\nabla_\eta A\|_\infty
                         +[\nabla_\eta A]_\alpha\le B.
 \label{fmg-smooth-assumption}
\end{equation}
The same assertion holds on an interior subcube containing
adjacent plaquettes of both parities.
\end{proposition}

\begin{proof}
For any \(A\) satisfying \eqref{fmg-smooth-assumption}, express
a \(\mu\nu\)-plaquette by the exact jet variables
\[
 a=A_\mu(x),\quad b=A_\nu(x),\quad
 c=\nabla_{\eta,\mu}A_\nu(x),\quad
 e=\nabla_{\eta,\nu}A_\mu(x).
\]
Its product is
\[
 P_\eta=
 e^{i\eta a}e^{i\eta b+i\eta^2c}
                  e^{-i\eta a-i\eta^2e}e^{-i\eta b}.
\]
The constant term is \(I\) and the first coefficient in \(\eta\)
vanishes. Consequently \((P_\eta-I)/\eta^2\) extends analytically
to \(\eta=0\), with uniformly bounded first jet derivatives
on each fixed bounded jet ball. Adjacent \(a,b\) differ by
at most \(B\eta\), and adjacent \(c,e\) differ by at most
\(B\eta^\alpha\). Hence
\begin{equation}
 \left\|(P_\eta(x)-I)/\eta^2
             -(P_\eta(x+e_\rho)-I)/\eta^2\right\|
                         \le C(B)\eta^\alpha .
 \label{fmg-curvature-holder}
\end{equation}
The same inequality holds for the difference of their norms.
But \(\eta^{-2}\|U(\partial p)-I\|\) is gauge invariant
at each plaquette. For \eqref{fmg-oscillating-curvature}
its adjacent values are zero and at least \(f/2\) for
sufficiently small \(\eta\). This contradicts
\eqref{fmg-curvature-holder}.
\end{proof}

This is an obstruction to every gauge, not merely to the chosen
tree gauge. The example is already in that tree gauge, where
\(A_1=\eta f x_2\chi(x_1)H\); its first derivative visibly grows
like \(f\ell/\eta\). The stronger invariant argument above
does not assert that a bounded \(C^1\) gauge alone is impossible.
It excludes the fixed physical \(C^{1,\alpha}\) background ball
needed here.

There is no contradiction with the complete-density method.
A raw fine field is not a uniformly smooth classical background.
The useful decomposition is a regular constrained background
plus rough fluctuations and explicitly retained exceptional
regions. A plaquette tail or a local logarithmic chart alone
does not construct that decomposition.

\subsection{Flat torus holonomies and transition functions}
\label{fmg-holonomy}

On a periodic direction of \(N\) bonds and physical circumference
\(\ell=N\eta\), consider
\[
       U_1(x)=e^{i\theta H/N},\qquad U_\mu(x)=I\quad(\mu\ne1).
\]
Its plaquettes are all \(I\), whereas the noncontractible loop is
\(e^{i\theta H}\), unchanged up to conjugacy by a periodic gauge.
If some periodic gauge has \(U_b^v=e^{i\eta A_b}\) with
\(\|A\|_\infty\le b\), telescoping the loop gives
\begin{equation}
                  \|e^{i\theta H}-I\|\le N\eta b=\ell b.
 \label{fmg-holonomy-bound}
\end{equation}
Thus zero plaquette curvature cannot force an arbitrarily small
global physical potential. Local tree gauges trivialize this
flat field on contractible cubes, but their transition data
retain the holonomy. On overlaps the exact relation is
\begin{equation}
 U^{(i)}(x,y)=h_{ij}(x)U^{(j)}(x,y)h_{ij}(y)^{-1},
 \qquad h_{ij}=v_i v_j^{-1}.
 \label{fmg-transitions}
\end{equation}
Those data must be retained or integrated consistently. They
cannot be set to identity because local plaquettes are small.

\subsection{A scale-comparison fact supplied by local shooting}
\label{fmg-couplings}

Our actual weak finite trajectories obey
\(\gamma_j=g_j^{-2}\ge\Gamma_{\mathrm w}>0\) and
\(|\gamma_{j+1}-\gamma_j|\le D_{\mathrm{abs}}\).
Telescoping yields, for any two stages,
\begin{equation}
 \frac{\gamma_i}{\gamma_j}
       \le1+\frac{D_{\mathrm{abs}}}{\Gamma_{\mathrm w}}|i-j|,
 \qquad
 \frac{g_i}{g_j}
       \le\left(1+\frac{D_{\mathrm{abs}}}
                          {\Gamma_{\mathrm w}}|i-j|\right)^{1/2}.
 \label{fmg-coupling-comparison}
\end{equation}
For the second inequality apply the first with \(i,j\) interchanged.
Both reverse ratios satisfy the same bounds. No one-step
monotonicity is used. This is a useful polynomial interscale
comparison; it is not asserted to match every source radius
inequality without checking its constants and domains.

\subsection{The next normalized interface}
\label{fmg-next}

A direct full-measure route needs an exact gauge-invariant
good/bad partition and a normalized conditional representation.
Its good component must be the actual stationary-plus-fluctuation
integral with its actual retained window. Its complementary
replacement must be summable in component size, scale age and
boundary geometry, uniformly in the allowed exterior data.
A fixed-box conditional Wilson estimate is a genuine first input;
its constants need not survive growing components or a nonlinear
average constraint. Conversely, a normalized estimate on a
coarse-average fiber cannot be inferred by conditioning an
unconditional tail without controlling that fiber's normalization.

The import route makes the same obligations through the published
complete-density class instead of a new direct proof. Neither
route has been completed by the local trajectory or by the
geometric lemmas in this section. No continuum-existence conclusion
is asserted here.


\section{The remaining construction theorem}
\label{fmb:remaining}

The progress now has a precise measure-theoretic meaning. For one
specified completion of the local coarse map, the full finite
positive Wilson pushforward is the sum of positive history measures.
At the bare level, large-curvature cells and explicit fixed-box
gauge/minimizer defects have uniform-volume joint probability
bounds. Neither assertion discards a bad sector or a vacuum factor.

What remains is to make the two constructions refer to the same
iterated local coordinates and to control their interaction. A
useful next theorem is a \emph{complete-density transfer theorem}
for the actual chosen coarse map. It would need to provide:
\begin{enumerate}
\item A spatial good/bad decomposition on determining sets with
collars, prescribed exterior data, and retained holonomy/transition
data. Its good-component integral must equal the existing
constrained stationary-plus-fluctuation integral, with its actual
independent windows, Haar factors, pivot factors and vacuum
normalization.
\item A normalized replacement on each complementary component.
The numerator and denominator must be bounded together, uniformly
in the admissible exterior field, component geometry, preceding
depth and scale history. A bound sufficient for summation over
components and their histories is required. Bare probability
rarity or a Gaussian conditional estimate does not by itself
provide this normalized replacement.
\item Preservation of a specified complete-density class under the
full step, with local analytic/source domains and coupling
extraction belonging to that same class. The weak local shooting
from Notebook 37 can be used only after its recursion has been
identified with this full-density recursion, or after a shooting
theorem has been proved for the latter.
\end{enumerate}
The published large-field induction discussed above is a possible
route to this theorem after its hypotheses have been matched.
The new direct probability estimates are another input to a proof;
they do not already verify the published induction class.

\paragraph{Why local rarity is not a global chart theorem.}
For any fixed family of \(M\) boxes, the one-box defect estimate
gives a union bound proportional to \(M\). If the number of boxes
diverges with the cutoff, a small one-box probability does not
imply that all boxes are good. For example, on the local tuned
trajectories \(g_{\rm bare}\) has order \(n^{-1/2}\), while a
logarithmic window yields a displayed estimate of order
\(\exp[-c(\log n)^{2a}]\). This bound does not defeat a number
of microscopic boxes growing exponentially in \(n\). This is
a limitation of that union-bound inference, not a lower bound
on the actual defect probability. The purpose of a component
expansion is to allow defects and sum them with their correct
weights, rather than assume a globally good field.

\paragraph{After the complete-density theorem.}
Continuum existence still requires common estimates for a
sufficient gauge-invariant observable algebra, with renormalized
short-distance insertions and the exact same denominator. Those
estimates must produce a common subsequential continuum hierarchy,
reflection positivity inherited from one positive regulator
sequence, full Euclidean covariance, and the regularity needed
for reconstruction. A quantitative lower witness and
identification of the resulting fields and dynamics are needed
to establish a nontrivial Yang--Mills theory. The physical mass
gap is a subsequent assertion. The Gaussian gaps used in the
local analysis are not that spectral statement.

No uniqueness of the continuum limit is imposed merely for the
sake of an existence proof. Nor is a particular contact-term
convention required when an admissible consistent choice would
suffice. The unresolved requirements above concern the actual
measure, observable bounds and physical identification.

\section*{Verification record and source access}

The accompanying research notes preserve the exact finite-measure
construction, the self-contained Wilson chessboard proof, the
fixed-box conditional proof and their independent internal audits.
The assembled source was checked for duplicate and missing labels
and compiled with the desktop LaTeX compiler. These checks establish
document consistency and typesetting; the mathematical evidence is
the arguments and their stated hypotheses.

Primary references used to locate the complete-density interface are
T.~Ba{\l}aban, \emph{Convergent renormalization expansions for lattice
gauge theories}, Communications in Mathematical Physics
\textbf{119} (1988), 243--285,
\href{https://doi.org/10.1007/BF01217741}{doi:10.1007/BF01217741},
and \emph{Large field renormalization. II. Localization,
exponentiation, and bounds for the \(R\) operation},
Communications in Mathematical Physics \textbf{122} (1989),
355--392,
\href{https://doi.org/10.1007/BF01238433}{doi:10.1007/BF01238433}.
The exact Section 2 induction class in the first paper, and the
weak-coupling hypothesis in Theorem 1 of the second, are retained.
No unchecked OCR constants or powers are supplied as theorem
hypotheses. The first large-field paper's detailed replacement
formula has not been newly verified from an accessible original
in this pass; its exact audit remains part of any import route.

\end{document}
